\begin{question}[ID=1.1, subtitle={Teorema di Lagrange per $f(x)=x^3+2x$ su $[-2,1]$}, tags={lagrange, teoremi}, difficulty=facile]
Verifica se la funzione \(f(x) = x^3+2x\) soddisfa le ipotesi del teorema di Lagrange nell'intervallo \([-2;\ 1]\) e, in caso affermativo, trova il punto \(c\) che soddisfa il teorema.
\begin{solution}
$f$ continua e derivabile su $[-2,1]$. $f'(c)=\dfrac{f(1)-f(-2)}{3}=\dfrac{3-(-12)}{3}=5$. $3c^2+2=5 \Rightarrow c^2=1 \Rightarrow c=-1$
\end{solution}
\end{question}

%-----------------------------%

\begin{question}[ID=1.2, subtitle={Teorema di Lagrange per $f(x)=3/\sqrt{x}$ su $[4,9]$}, tags={lagrange, teoremi}, difficulty=media]
Verifica se la funzione \(f(x) = \dfrac{3}{\sqrt{x}}\) soddisfa le ipotesi del teorema di Lagrange nell'intervallo \([4;\ 9]\) e, in caso affermativo, trova il punto \(c\) che soddisfa il teorema.
\begin{solution}
$f'(c)=\dfrac{f(9)-f(4)}{5}=\dfrac{1-\tfrac{3}{2}}{5}=-\dfrac{1}{10}$. $-\dfrac{3}{2c^{3/2}}=-\dfrac{1}{10} \Rightarrow c^{3/2}=15 \Rightarrow c=15^{2/3}$
\end{solution}
\end{question}

%-----------------------------%

\begin{question}[ID=1.3, subtitle={Teorema di Lagrange non applicabile: funzione a tratti su $[0,2]$}, tags={lagrange, teoremi, non applicabile}, difficulty=media]
Data la funzione
\[
f(x) = \begin{cases}
x^2 + 1, & \text{se } x \leq 1 \\
3 - x, & \text{se } x > 1
\end{cases}
\]
verifica se soddisfa le ipotesi del teorema di Lagrange nell'intervallo \([0;\ 2]\) e, in caso affermativo, trova il punto \(c\) che soddisfa il teorema.
\begin{solution}
Non applicabile: $f$ è continua in $x=1$ ma non derivabile ($f'(1^-)=2 \ne f'(1^+)=-1$).
\end{solution}
\end{question}

%-----------------------------%

\begin{question}[ID=1.4, subtitle={Teorema di Lagrange per funzione a tratti su $[0,2]$}, tags={lagrange, teoremi}, difficulty=media]
Data la funzione
\[
f(x) = \begin{cases}
x^2 -4x+5, & \text{se } x \leq 2 \\
x, & \text{se } x > 2
\end{cases}
\]
verifica se soddisfa le ipotesi del teorema di Lagrange nell'intervallo \([0;\ 2]\) e, in caso affermativo, trova il punto \(c\) che soddisfa il teorema.
\begin{solution}
Su $[0,2]$ vale $f(x)=x^2-4x+5$; continua e derivabile. $f'(c)=\dfrac{f(2)-f(0)}{2}=\dfrac{1-5}{2}=-2$. $2c-4=-2 \Rightarrow c=1$
\end{solution}
\end{question}

%-----------------------------%

\begin{question}[ID=1.5, subtitle={Teorema di Lagrange per $f(x)=2e^x+x$ su $[0,1]$}, tags={lagrange, teoremi}, difficulty=media]
Verifica se la funzione \(f(x) = 2e^x + x\) soddisfa le ipotesi del teorema di Lagrange nell'intervallo \([0;\ 1]\) e, in caso affermativo, trova il punto \(c\) che soddisfa il teorema.
\begin{solution}
$f'(c)=\dfrac{f(1)-f(0)}{1}=(2e+1)-2=2e-1$. $2e^c+1=2e-1 \Rightarrow e^c=e-1 \Rightarrow c=\ln(e-1)$
\end{solution}
\end{question}

%-----------------------------%

\begin{question}[ID=1.6, subtitle={Teorema di Lagrange non applicabile: $f(x)=4/x$ su $[-1/2,1]$}, tags={lagrange, teoremi, non applicabile}, difficulty=facile]
Verifica se la funzione \(f(x) = \dfrac{4}{x}\) soddisfa le ipotesi del teorema di Lagrange nell'intervallo \(\bigl[-\tfrac{1}{2};\ 1\bigr]\) e, in caso affermativo, trova il punto \(c\) che soddisfa il teorema.
\begin{solution}
Non applicabile: $f$ è discontinua in $x=0 \in \bigl[-\tfrac{1}{2}, 1\bigr]$.
\end{solution}
\end{question}

%-----------------------------%

\begin{question}[ID=2.1, subtitle={Parametro $a$ per Lagrange: $x^3+ax$ su $[-1,1]$}, tags={lagrange, teoremi, parametrico}, difficulty=difficile]
Trova il valore di \(a\) tale che la funzione \(f(x) = x^3 + ax\) soddisfi il teorema di Lagrange nell'intervallo \([-1;\ 1]\) e determina il punto \(c\) che soddisfa il teorema.
\begin{solution}
$f$ è continua e derivabile per ogni $a$. $f'(c)=\dfrac{f(1)-f(-1)}{2}=\dfrac{(1+a)-(-1-a)}{2}=1+a$. $3c^2+a=1+a \Rightarrow c^2=\dfrac{1}{3} \Rightarrow c=\pm\dfrac{1}{\sqrt{3}}$ (per ogni $a$).
\end{solution}
\end{question}

%-----------------------------%

\begin{question}[ID=2.2, subtitle={Parametro $b$ per Lagrange: $1/x$ su $[1,b]$ con $c=2$}, tags={lagrange, teoremi, parametrico}, difficulty=media]
Data la funzione \(f(x) = \dfrac{1}{x}\) nell'intervallo \([1;\ b]\), determina il valore di \(b > 1\) tale che il punto \(c\) che soddisfa il teorema di Lagrange sia \(c = 2\).
\begin{solution}
$f'(c)=-\dfrac{1}{c^2}=-\dfrac{1}{4}$. $\dfrac{f(b)-f(1)}{b-1}=\dfrac{\tfrac{1}{b}-1}{b-1}=-\dfrac{1}{b}=-\dfrac{1}{4} \Rightarrow b=4$
\end{solution}
\end{question}

%-----------------------------%

\begin{question}[ID=2.3, subtitle={Parametri $a,b$ per Lagrange: funzione a tratti su $[-2,0]$}, tags={lagrange, teoremi, parametrico}, difficulty=difficile]
Trova il valore di \(a\) e \(b\) tale che la funzione
\[
f(x) = \begin{cases}
x^2 +ax-1, & \text{se } x < -1 \\
-4x^2+x+b, & \text{se } x \ge -1
\end{cases}
\]
soddisfi il teorema di Lagrange nell'intervallo \([-2;\ 0]\) e determina il punto \(c\) che soddisfa il teorema.
\begin{solution}
Continuità in $x=-1$: $1-a-1=-4-1+b \Rightarrow b=5-a$. Derivabilità: $-2+a=8+1=9 \Rightarrow a=11,\ b=-6$. $c_1=-\tfrac{9}{4}$ oppure $c_2=-\tfrac{11}{16}$.
\end{solution}
\end{question}

%-----------------------------%

\begin{question}[ID=2.4, subtitle={Lagrange per $|x^2-6x+5|$ su $[0,2]$ e $[2,a]$}, tags={lagrange, teoremi, valore assoluto}, difficulty=difficile]
Data la funzione $f(x) = |x^2 - 6x + 5|$, determina se soddisfa il teorema di Lagrange nell'intervallo \([0;\ 2]\). Determina inoltre il valore massimo del parametro $a$ affinché il teorema di Lagrange sia valido nell'intervallo $[2, a]$ e, in caso affermativo, determina i valori che soddisfano il teorema.
\begin{solution}
Su $[0,2]$: non applicabile ($f$ non derivabile in $x=1$). Su $[2,a]$ con $a \le 5$: $f(x)=-x^2+6x-5$, $c=1+\dfrac{a}{2} \in (2, a)$. Valore massimo $a=5$.
\end{solution}
\end{question}

%-----------------------------%

\begin{question}[ID=2.5, subtitle={Lagrange per $|x^2-4x+3|$ su $[0,2]$ e $[2,a]$}, tags={lagrange, teoremi, valore assoluto}, difficulty=difficile]
Data la funzione $f(x) = |x^2 - 4x + 3|$, determina se soddisfa il teorema di Lagrange nell'intervallo \([0;\ 2]\). Determina inoltre il valore massimo del parametro $a$ affinché il teorema di Lagrange sia valido nell'intervallo $[2, a]$ e, in caso affermativo, determina i valori che soddisfano il teorema.
\begin{solution}
Su $[0,2]$: non applicabile ($f$ non derivabile in $x=1$). Su $[2,a]$ con $a \le 3$: $f(x)=-x^2+4x-3$, $c=1+\dfrac{a}{2} \in (2, a)$. Valore massimo $a=3$.
\end{solution}
\end{question}
